\chapter*{Notation and conventions}
\addcontentsline{toc}{chapter}{Notation and conventions}
% STATUS: written 2026-11-06 (records the notations the chapters actually use;
% each symbol is defined in its chapter, and this list is the index to them.
% Operators come from style/encyclopaedia.sty; no new notation is introduced here.)
\label{page:notation}

The book keeps one notation throughout.
A symbol is defined where it first appears, and the chapters reuse it without
giving it a second meaning.
This page lists the symbols a reader meets most often, so that a chapter can be
opened at any point.

\section*{Rings and algebras}

$R$ and $A$ are a ring and an algebra; the book says which when the distinction
matters. $\mathbb{F}$ is a field, and $k$ is a field when the work is
geometric. $\operatorname{char}R$ is the characteristic of $R$, and $Z(R)$ is
its centre.
An algebra is associative and unital unless a chapter says otherwise;
Chapter~\ref{ch:nonassociative} drops associativity and says so at each step.

$M_n(R)$ is the ring of $n\times n$ matrices over $R$.
A triangular algebra is written $A=\operatorname{Tri}(A_1,A_2,M)$ with $M$ the
bimodule between the two parts.
$R[x]$ and $R[x_1,\dots,x_n]$ are polynomial rings, and $k(x)$ is the field of
fractions when $k$ is a domain.

\section*{Derivations and their relatives}

$\operatorname{Der}(A)$ is the set of derivations of $A$, an additive group and
a Lie algebra under the commutator.
$\operatorname{Inn}(A)$ is the set of inner derivations, those of the form
$x\mapsto[\,a,x\,]$, so that
$\operatorname{HH}^{1}(A,A)\cong\operatorname{Der}(A)/\operatorname{Inn}(A)$ for
the classes where that identification holds.

$\operatorname{JDer}(A)$ and $\operatorname{LDer}(A)$ are the Jordan and the Lie
derivations. $\operatorname{ad}_a$ is the inner derivation $x\mapsto[a,x]$.
A derivation is \emph{locally nilpotent} when some power of it kills each
element, and the set of these is written $\operatorname{LND}(A)$. % ours: notation-page term, defined in Chapter 16

$\operatorname{HH}^{\bullet}(A,M)$ is the Hochschild cohomology of $A$ with
coefficients in the bimodule $M$, and $\operatorname{HH}^{\bullet}(A)$ is the
case $M=A$.
Low degrees are written out where they are used: the first cohomology group is
the derivations modulo the inner ones.

\section*{Maps and products}

$[x,y]=xy-yx$ is the commutator and $x\circ y=xy+yx$ is the anti-commutator.
A map is \emph{additive} when it respects addition, with no linearity assumed
over a ring; the book states linearity separately wherever a theorem needs it. % ours: notation-page convention

For maps with a parameter, $\sigma$ and $\tau$ are endomorphisms of the algebra,
and $\alpha$, $\beta$ are scalars or endomorphisms as the chapter says.
A $\sigma$-derivation satisfies $d(xy)=d(x)y+\sigma(x)d(y)$, and the book gives
the defining identity for each twisted kind at the point where it is named.

\section*{Continuity and operators}

For a Banach algebra, $\lVert\cdot\rVert$ is the norm, and a derivation is
\emph{bounded} when it is bounded as a linear map. % ours: notation-page term, defined in Chapter 18
The book asks, for each class, whether continuity has to be assumed or comes for
free; the question is fixed in
Section~\ref{sec:basic-theory:question-families} and answered by class in the
chapters.

$\mathcal{B}(H)$ is the algebra of bounded operators on a Hilbert space $H$, and
$\mathcal{K}(H)$ is the ideal of compact operators.
A von Neumann algebra is written $M$.

\section*{Conventions}

A theorem recalled from a source carries the source in its header, with the page
or the numbered item.
A result proved in this book is marked as one of its own.
A citation after a sentence gives the page on which the cited source says what
the sentence attributes to it.

The book uses British spelling in the text and keeps the spelling of a quoted
title.
Section and theorem numbers are stable within an edition; a later edition may
add a section at the end of a chapter, which leaves earlier numbers alone.
